\documentclass{article} % say
\usepackage{amsmath}
\usepackage{amssymb}
\usepackage{gensymb}
\usepackage{hyperref}
\usepackage{tikz}
\usepackage{tkz-euclide}
\usetikzlibrary {babel,angles,quotes,positioning}
\usepackage[utf8]{inputenc}
\usepackage[russian]{babel}

\usepackage{mathtools}
\DeclarePairedDelimiter\ceil{\lceil}{\rceil}
\DeclarePairedDelimiter\floor{\lfloor}{\rfloor}
\begin{document}

\textbf{Задача}: Найти такой набор натуральных чисел, что его сумма равна $1976$, а произведение максимально.

\clearpage

\textbf{Решение}: 

Для начала докажем, что среди всех наборов, которые подходят под решение задачи, существует такой, что максимальное число в нём не превосходит $3$. Допустим это не так, выберем из всех подходящих наборов такой, что максимальное число в нём (назовём его $n$) минимально. Заменим это число на два числа: $\floor{\frac{n}{2}}$ и $\ceil{\frac{n}{2}}$. Заметим, что сумма чисел в наборе не поменялась, докажем, что произведение не уменьшилась.

Пусть $n$ чётно. Тогда необходимо доказать, что: $\frac{n^2}{4} \geqslant n$, что для натуральных $n$ выполняется, если $n \geqslant 4$.

Пусть $n$ нечётно. Тогда из $\frac{n - 1}{2} \cdot \frac{n + 1}{2} \geqslant n$, следует, что $n^2 - 1 \geqslant 4n$, что верно для нечётных $n > 3$.

Таким образом существует набор вида $a$ единиц, $b$ двоек, $c$ троек, удовлетворяющий условиям задачи. Будем задавать набор тройкой $(a, b, c)$, причём тройку будем называть хорошей, если $a + 2b + 3c = 1976$, т.е. сумма чисел из набора равна $1976$. Очевидно, что произведение чисел в таком наборе равно $2^b \cdot 3^c$.

Докажем, что в оптимальном наборе $a < 2$. Действительно, если это не так, то мы можем из хорошей тройки $(a, b, c)$ получить хорошую тройку $(a - 2k, b + k, c)$, где $k = \floor{\frac{a}{2}}$, так как сумма набора не изменится, а вот произведение возрастёт в $2^k$ раз.

Также, в оптимальном наборе $b < 3$. Действительно, если это не так, то мы можем из хорошей тройки $(a, b, c)$ получить хорошую тройку $(a, b - 3k, c + 2k)$, где $k = \floor{\frac{b}{3}}$, так как сумма набора не изменится, а вот произведение возрастёт в $\left( \frac{9}{8} \right)^k$ раз.

Кроме того, в оптимальном наборе $a$ и $b$ не могут быть оба больше нуля. Действительно, если это не так, то мы можем из хорошей тройки $(a, b, c)$ получить хорошую тройку $(a - k, b - k, c + k)$, где $k = \min(a,b)$, так как сумма набора не изменится, а вот произведение возрастёт в $\left( \frac{3}{2} \right)^k$ раз.

Таким образом, искомый оптимальный набор можно найти среди следующих троек:

\begin{enumerate}
  \item $(0, 0, \frac{1976}{3})$,
  \item $(1, 0, \frac{1975}{3})$,
  \item \label{item:correct} $(0, 1, \frac{1974}{3})$,
  \item $(0, 2, \frac{1972}{3})$.
\end{enumerate}

Заметим, что только вариант \ref{item:correct} может быть ответом, так как только в нём $c$ является целым числом. Таким образом, ответ: одна двойка и $658$ троек.  

\textbf{Источник}: \url{https://youtu.be/o-UF8JlqZ_c?si=ZtiHRf2mMf1lhalB}

\end{document}


